\section{Objetivos}

\subsection{Objetivo General}
Desarrollar un prototipo de visualizaci\'on Web 3D interactiva basado en Unity WebGL, enfocado en la exploraci\'on t\'ecnica y el an\'alisis visual de hardware de alto rendimiento, optimizado mediante criterios de \textit{Technical Art} para su despliegue eficiente en navegadores web est\'andar.

\subsection{Objetivos Espec\'ificos}

Dise\~nar un \textit{pipeline} de optimizaci\'on de activos 3D que reduzca la carga poligonal de modelos CAD originales hasta un presupuesto geom\'etrico compatible con WebGL, manteniendo fidelidad visual mediante retopolog\'ia, limpieza de mallas y \textit{baking} de texturas.

Integrar en Unity URP un sistema de materiales y modos de visualizaci\'on que combine renderizado realista con herramientas de inspecci\'on t\'ecnica, procurando un \textit{frame time} igual o inferior a 33.33 ms (30 FPS o m\'as) en el dispositivo objetivo.

Programar un sistema de interacci\'on en C\# compilado a WebAssembly que permita navegaci\'on orbital, selecci\'on de piezas, consulta contextual de informaci\'on y uso de herramientas como vista explosionada, corte transversal y modos anal\'iticos del prototipo final.

Validar el rendimiento, la usabilidad y la carga cognitiva del prototipo mediante herramientas de perfilado, pruebas con usuarios, cuestionarios SUS, NASA-TLX Raw y protocolo \textit{Think-Aloud}, contrastando los resultados frente a las metas operativas definidas para el proyecto.

\newpage
